% Originally: content/week-03.tex, \subsection{Contraction mappings and the Banach fixed point theorem}

% Sonnet 5 (Medium)
\section{Contraction mappings and the Banach fixed point theorem}
\label{sec:banach_fixed_point}

% Not in Corsin's own notes either, but needed to justify the appeal to the "Banach Fixed Point
% Theorem" in the solution to ex:ai_contraction_cos below, and a standard companion to
% completeness (def:complete_metric_space), which is why it sits in this chapter.
% Transition: Claude Opus 5 (High)
Completeness guarantees that certain limits exist, but on its own it does not say what they are.
The following theorem is the standard way of turning that guarantee into a construction: it
identifies a single point, proves it exists, proves it is unique, and supplies an algorithm that
converges to it. Almost every existence proof in the second half of this course goes through it.

\begin{definition}[Contraction mapping]
\label{def:contraction_mapping}
Let $(X,d)$ be a metric space. A map $f : X \to X$ is a \newterm{contraction}
\germanfor{Kontraktion} if it is Lipschitz continuous (\cref{def:lipschitz_continuous}) with
some Lipschitz constant $L \in [0,1)$, i.e.,
\[d(f(x),f(y)) \leq L \cdot d(x,y) \quad \text{for all } x,y \in X.\]
\end{definition}

\begin{theorem}[Banach fixed point theorem]
\label{thm:banach_fixed_point}
Let $(X,d)$ be a nonempty complete metric space (\cref{def:complete_metric_space}) and let
$f : X \to X$ be a contraction (\cref{def:contraction_mapping}). Then $f$ has a unique
\newterm{fixed point} \germanfor{Fixpunkt} $x^* \in X$, i.e., $f(x^*) = x^*$; moreover, for any
$x_0 \in X$, the iterates $x_{n+1} := f(x_n)$ converge to $x^*$.
\end{theorem}

% Supplement: Sascha Brack/Class Notes/Week_02_Notes_Monday_Updated.pdf, p. 13
\begin{proof}[Proof idea]
The core idea is to create a Cauchy sequence by repeatedly applying the contraction mapping $f$.
Fix an arbitrary starting point $x_0 \in X$, and define the sequence by $x_{n+1} := f(x_n)$ for all $n \geq 0$.
Comparing consecutive terms gives:
\[
d(x_{n+1}, x_n) = d(f(x_n), f(x_{n-1})) \leq \lambda \cdot d(x_n, x_{n-1}).
\]
Iterating this inequality yields $d(x_{n+1}, x_n) \leq \lambda^n d(x_1, x_0)$. Since $\lambda < 1$, one can use the triangle inequality and the geometric series to show that the distance between any two terms $x_m$ and $x_n$ (for $m > n$) becomes arbitrarily small as $n \to \infty$. Thus, the sequence is Cauchy.
Because $X$ is a complete metric space, this sequence converges to some limit $x^* \in X$. Since contraction mappings are continuous, one can pass the limit inside the function to conclude $f(x^*) = x^*$.
\end{proof}

% Generator: Claude Opus 5 (High) --- Sascha's proof idea establishes existence only; uniqueness
% is two lines and is half of what the theorem claims.
\begin{proof}[Uniqueness of the fixed point]
Suppose $f(x) = x$ and $f(y) = y$. Then
\[d(x,y) = d\big(f(x),f(y)\big) \leq L\,d(x,y),\]
so $(1-L)\,d(x,y) \leq 0$. Since $L<1$ we have $1-L>0$, forcing $d(x,y) \leq 0$ and hence
$d(x,y)=0$, i.e.\ $x=y$.
\end{proof}

% Generator: Claude Opus 5 (High) --- FIG-06-01, cobweb diagrams for the fixed point iteration.
% This chapter had no figure at all, and the theorem's proof is a construction (iterate f from
% an arbitrary starting point) whose whole content is visual in one variable.
% Arithmetic check. LEFT: f(x) = 0.45x + 0.35, a contraction with L = 0.45 < 1. Fixed point
% x* = 0.35/(1-0.45) = 0.35/0.55 = 0.6364. From x0 = 0.05 the iterates are
%   0.05 -> 0.3725 -> 0.51763 -> 0.58293 -> 0.61232 -> 0.62554, climbing to x*.
% RIGHT: f(x) = x + 0.12, slope exactly 1, so L = 1 and NOT a contraction. The graph is
% parallel to the diagonal and never meets it, so there is no fixed point at all. From
% x0 = 0.10 the iterates are 0.22, 0.34, 0.46, 0.58, 0.70, 0.82, marching off to the right.
% This is the second bullet of rem:banach_hypotheses (the T(x) = x+1 example) drawn on [0,1].
\begin{center}
\begin{tikzpicture}[x={(3.4cm,0)}, y={(0,3.4cm)}, >=Latex,
    web/.style={BrickRed, thick},
    fgraph/.style={blue!65!black, very thick}]

  % ---------------- left: a genuine contraction ----------------
  \begin{scope}
    \draw[->, gray!70] (0,0) -- (1.1,0) node[right, font=\footnotesize, text=black] {$x$};
    \draw[->, gray!70] (0,0) -- (0,1.1);
    \draw[gray!55, dashed] (0,0) -- (1,1) node[above left, font=\scriptsize, text=black] {$y=x$};
    \draw[fgraph] (0,0.35) -- (1,0.80) node[right, font=\scriptsize] {$f$};

    \draw[web] (0.05,0.05) -- (0.05,0.3725) -- (0.3725,0.3725) -- (0.3725,0.51763)
      -- (0.51763,0.51763) -- (0.51763,0.58293) -- (0.58293,0.58293) -- (0.58293,0.61232)
      -- (0.61232,0.61232) -- (0.61232,0.62554);

    \fill[black] (0.6364,0.6364) circle (1.7pt);
    \node[font=\scriptsize, anchor=west] at (0.68,0.60) {$x^{*}$};
    \draw[gray!70] (0.05,0.02) -- (0.05,-0.02) node[below, font=\scriptsize, text=black] {$x_0$};
    \node[font=\scriptsize, align=center, anchor=north] at (0.5,-0.15)
      {contraction, $L=0.45$\\ iterates climb to the unique $x^{*}$};
  \end{scope}

  % ---------------- right: slope exactly 1, no fixed point ----------------
  \begin{scope}[shift={(1.42,0)}]
    \draw[->, gray!70] (0,0) -- (1.1,0) node[right, font=\footnotesize, text=black] {$x$};
    \draw[->, gray!70] (0,0) -- (0,1.1);
    \draw[gray!55, dashed] (0,0) -- (1,1) node[below right, font=\scriptsize, text=black] {$y=x$};
    \draw[fgraph] (0,0.12) -- (0.88,1.0) node[above left, font=\scriptsize] {$f$};

    \draw[web] (0.10,0.10) -- (0.10,0.22) -- (0.22,0.22) -- (0.22,0.34) -- (0.34,0.34)
      -- (0.34,0.46) -- (0.46,0.46) -- (0.46,0.58) -- (0.58,0.58) -- (0.58,0.70)
      -- (0.70,0.70) -- (0.70,0.82) -- (0.82,0.82) -- (0.82,0.94);
    \draw[web, ->] (0.82,0.94) -- (0.94,0.94);

    \draw[gray!70] (0.10,0.02) -- (0.10,-0.02) node[below, font=\scriptsize, text=black] {$x_0$};
    \node[font=\scriptsize, align=center, anchor=north] at (0.5,-0.15)
      {$L=1$: graph parallel to $y=x$\\ no fixed point, iterates escape};
  \end{scope}
\end{tikzpicture}
\end{center}

The left panel is the proof in miniature. Each step moves vertically to the graph of $f$ and then
horizontally back to the diagonal, and because the graph is shallower than the diagonal, the steps
shrink by a factor of at least $L$ every time. That is precisely the geometric-series estimate the
proof runs on. The right panel shows what a single missing hypothesis costs: with slope exactly
$1$ the graph never meets the diagonal, so no fixed point exists and the same iteration marches
off instead of converging.

\begin{remark}[Where each hypothesis is used]
\label{rem:banach_hypotheses}
Both hypotheses of \cref{thm:banach_fixed_point} are indispensable, and each fails in a
different way.
\begin{itemize}
  \item \textbf{Completeness} is what supplies the limit. On the incomplete space
    $X := (0,1]$, the map $f(x) := x/2$ is a contraction with $L=\tfrac12$ and has no fixed
    point in $X$, since the iterates march towards $0$, which is exactly the point missing.
  \item \textbf{$L<1$ strictly} is what forces the iterates together. The map
    $f(x) := x + \tfrac{1}{1+e^{x}}$ on the complete space $\mathbb{R}$ satisfies the strict
    inequality $|f(x)-f(y)| < |x-y|$ for $x \neq y$, yet has no fixed point, since
    $f(x)>x$ everywhere. A supremum of $1$ that is never attained is still not good enough.
\end{itemize}
The uniqueness proof above shows the second point cleanly: the whole argument is the division by
$1-L$, which is exactly what $L=1$ forbids.
\end{remark}

% Supplement: Analysis_II_Script_v1.pdf, p. 20
\begin{exercise}[Both hypotheses of Banach's theorem are sharp]
\label{ex:banach_hypotheses_sharp}
Find examples for:
\begin{enumerate}[label=\textbf{(\alph*)}]
  \item A Lipschitz contraction $T : X \to X$ on a non-complete metric space $X$ that has no
    fixed point.
  \item A complete metric space $(X,d)$ and an \newterm{isometry} (a mapping $T : X \to X$ with
    $d(T(x_1),T(x_2)) = d(x_1,x_2)$ for all $x_1,x_2$) that has no fixed point.
\end{enumerate}
\end{exercise}

\begin{remark}
Part \textbf{(b)} sharpens \cref{rem:banach_hypotheses}'s second bullet: that example already
showed a Lipschitz constant of exactly $1$ can defeat the theorem, but its map $T$ was not an
isometry ($|T(x)-T(y)| < |x-y|$ strictly, just with no uniform gap). An isometry is the most
rigid failure imaginable, not shrinking distances even a little, and it \emph{still} need not have
a fixed point, even on a complete space. Note that unboundedness is doing real work in the
counterexample: on a \emph{compact} space an isometry into itself is automatically surjective, and
even that is not enough to force a fixed point. An irrational rotation of the circle $S^1$, which
is compact and an isometry, has none either. Fixed points of isometries are a
genuinely separate question from \cref{thm:banach_fixed_point}, not one it settles as a special
case.
\end{remark}

\begin{aiexercise}[{Contraction Mapping on $[0,1]$}]
\label{ex:ai_contraction_cos}
% Generator: Gemini 3.6 Flash (Medium)
Let $f \colon [0,1] \to [0,1]$ be defined by $f(x) := \frac{1}{2}\cos(x)$.
\begin{enumerate}[label=\textbf{(\alph*)}]
  \item Show that $f$ is a contraction mapping (\cref{def:contraction_mapping}) with Lipschitz
    constant $L = \frac{1}{2}\sin(1) < 1$.
  \item Deduce, using \cref{thm:banach_fixed_point}, that $f$ has a unique fixed point
    $x^* \in [0,1]$.
\end{enumerate}
\end{aiexercise}

% Supplement: Linus Lüchinger/Slides/Slides-03-02.pdf, slide 16;
%             Linus Lüchinger/Slides/Slides-03-05.pdf, slide 3
\begin{exercise}[Does the fixed-point theorem apply?]
\label{ex:bfpt_applicability_table}
\cref{rem:banach_hypotheses} shows that each hypothesis is needed by exhibiting one failure
apiece. Linus drills the same point the other way round, as a checklist: for each row below,
decide whether \cref{thm:banach_fixed_point} \emph{applies} to $T : X \to X$, and separately
determine the set of fixed points of $T$ in $X$. The two questions have different answers more
often than one expects.
\begin{center}
\begin{tabular}{lll}
\textbf{$T(x)$} & \textbf{$X$} & \textbf{Applies? / fixed points} \\
\hline
$\sin(x)$ & $\mathbb{R}$ & \\
$\tfrac{9}{10}\cos(x)$ & $\mathbb{R}$ & \\
$x^2$ & $\mathbb{R}$ & \\
$\tfrac12x^2$ & $(-1,1)$ & \\
$\tfrac12\big(x + \tfrac2x\big)$ & $\mathbb{Q}\cap(1,2)$ &
\end{tabular}
\end{center}
\end{exercise}


% --- exercise relocated here from the dissolved sheet 3 ---

% Originally: content/exercise-sheets/week-03.tex, problem 3.5
% Source: exercises/Ex3_Analysis2_eng.pdf

\begin{exercise}[Multiple choice (fixed points)]
\label{ex:3.5}
Select all the statements below that are true.
\begin{enumerate}[label=\textbf{(\alph*)}]
  \item If you spread a map of Zurich on your desk, then one point of the desk will coincide with
    its representation on the map (in an ideal world).
  \item If $f : [0,1] \to [0,2]$ is $\tfrac{1}{2}$-Lipschitz, then $f$ has a fixed point.
  \item If $f : [0,2] \to [0,1]$ is $\tfrac{1}{2}$-Lipschitz, then $f$ has a fixed point.
  \item If $f : [0,1] \cup [2,3] \to [0,1] \cup [2,3]$ is continuously differentiable with
    $|f'(x)| < 1$ for all $x \in [0,1]\cup[2,3]$, then it has a fixed point.
  \item \textbf{(*)} If $f : [0,1] \to [0,1]$ is differentiable with $|f'(x)| < 1$ for all
    $x \in (0,1)$, then it has a fixed point.
\end{enumerate}
\exinfo{This exercise is Problem 3.5 of Problem Sheet 3, Analysis II, Spring Semester 2026.}
\end{exercise}
